\section{Arquitectura de Software}

El sistema a implementar para cumplir con los requisitos funcionales pueden ser
descripto en función de los siguientes:

• Ser accesible a múltiples usuarios.
• Permitir que los usuarios puedan interactuar con el sistema cumpliendo diferentes roles y
garantizar que ciertas funcionalidades estén disponibles sólo para algunos roles.
• Realizar validaciones sobre las operaciones que los usuarios intentan llevar a cabo, de
manera de garantizar el cumplimiento de las reglas de negocio y la consistencia de la
información almacenada.
• Permitir operaciones tanto de actualización como de consulta, de lo cual se deriva que se
debe almacenar el resultado de las operaciones efectuadas en algún medio persistente. 

Se decidió utilizar un patrón de arquitectura Enterprise, vinculado con el patron de layer. Se definieron las siguientes capas:
\begin{itemize}
        \item \textbf{Capa de presentación}: Encargada en presentar la información a los distintos usuario y permitiéndo, a partir de eventos generados por el usuario, interactuar con la capa de servicio.
        \item \textbf{Capa de servicio}: Brinda las funcionalidades necesarias para los distintos casos de uso a la capa de presentación e interactua con la capa de negocio.  
        \item \textbf{Capa de negocio}: Implementa todas las reglas de negocio necesarias para que se cumplan los requisitos pedidos.
        \item \textbf{Capa de persistencia}: Ofrece a la capa de negocio el servicio de persistir, almacenar y recuperar la información.
\end{itemize}